% GENERATED FILE. Edit canonical lessons, then run npm run generate:books.
% canonical-source-hash: fnv1a64:96947d5321f63c16
% canonical-lessons: SA-C198-harinah, SA-C198-citrakah, SA-C198-gandakah, SA-C198-nidam, SA-C198-andam

\chapter{Deer, Leopard, Rhinoceros, Nest, Egg}
\label{ch:sa-deer-leopard-rhinoceros-nest-egg}
\hlchaptermodality{198}

\begin{chapteropening}
\textbf{By the end of this chapter:} \emph{I can say a deer, a leopard, a rhinoceros, a nest and an egg.}

Complete the last lesson of chapter 198: I can say a deer, a leopard, a rhinoceros, a nest and an egg.
\end{chapteropening}

\section[\texorpdfstring{\sk{हरिणः} (hariṇaḥ)}{हरिणः (hariṇaḥ)}]{\sk{हरिणः} (hariṇaḥ) — a deer}

\label{lesson:SA-C198-harinah}



\begin{quote}
Before the new one: say the Sanskrit for dew, then the Sanskrit for moonlight.
\end{quote}

\subsection*{You'll want to know: \sk{हरिणः}}
\textbf{\sk{हरिणः}} — \emph{hariṇaḥ} — \textquotedblleft{}a deer\textquotedblright{}.

\begin{grammarlens}[title={What you've built}]
One more word for this chapter.
\end{grammarlens}

\subsection*{Your turn}
\begin{itemize}
\raggedright
  \item \emph{Say it:} \emph{hariṇaḥ}
  \item \emph{Say it:} \emph{hariṇaḥ}, once more
  \item \emph{Say it:} say \emph{hariṇaḥ}
  \item \emph{Recall:} say the Sanskrit for dew, then the Sanskrit for moonlight, then say \emph{hariṇaḥ} again
\end{itemize}

\begin{tcolorbox}[breakable,colback=teal!4,colframe=teal!35!black,title={Before you move on}]
What does \sk{हरिणः} mean? (\textquotedblleft{}A deer\textquotedblright{}.) Say it once more.
\end{tcolorbox}
\section[\texorpdfstring{\sk{चित्रकः} (citrakaḥ)}{चित्रकः (citrakaḥ)}]{\sk{चित्रकः} (citrakaḥ) — a leopard}

\label{lesson:SA-C198-citrakah}



\begin{quote}
Before the new one: say the Sanskrit for moonlight, then the Sanskrit for a deer.
\end{quote}

\subsection*{You'll want to know: \sk{चित्रकः}}
\textbf{\sk{चित्रकः}} — \emph{citrakaḥ} — \textquotedblleft{}a leopard\textquotedblright{}.

\begin{grammarlens}[title={What you've built}]
One more word for this chapter.
\end{grammarlens}

\subsection*{Your turn}
\begin{itemize}
\raggedright
  \item \emph{Say it:} \emph{citrakaḥ}
  \item \emph{Say it:} \emph{citrakaḥ}, once more
  \item \emph{Say it:} say \emph{citrakaḥ}
  \item \emph{Recall:} say the Sanskrit for moonlight, then the Sanskrit for a deer, then say \emph{citrakaḥ} again
\end{itemize}

\begin{tcolorbox}[breakable,colback=teal!4,colframe=teal!35!black,title={Before you move on}]
What does \sk{चित्रकः} mean? (\textquotedblleft{}A leopard\textquotedblright{}.) Say it once more.
\end{tcolorbox}
\section[\texorpdfstring{\sk{गण्डकः} (gaṇḍakaḥ)}{गण्डकः (gaṇḍakaḥ)}]{\sk{गण्डकः} (gaṇḍakaḥ) — a rhinoceros}

\label{lesson:SA-C198-gandakah}



\begin{quote}
Before the new one: say the Sanskrit for a deer, then the Sanskrit for a leopard.
\end{quote}

\subsection*{You'll want to know: \sk{गण्डकः}}
\textbf{\sk{गण्डकः}} — \emph{gaṇḍakaḥ} — \textquotedblleft{}a rhinoceros\textquotedblright{}.

\begin{grammarlens}[title={What you've built}]
One more word for this chapter.
\end{grammarlens}

\subsection*{Your turn}
\begin{itemize}
\raggedright
  \item \emph{Say it:} \emph{gaṇḍakaḥ}
  \item \emph{Say it:} \emph{gaṇḍakaḥ}, once more
  \item \emph{Say it:} say \emph{gaṇḍakaḥ}
  \item \emph{Recall:} say the Sanskrit for a deer, then the Sanskrit for a leopard, then say \emph{gaṇḍakaḥ} again
\end{itemize}

\begin{tcolorbox}[breakable,colback=teal!4,colframe=teal!35!black,title={Before you move on}]
What does \sk{गण्डकः} mean? (\textquotedblleft{}A rhinoceros\textquotedblright{}.) Say it once more.
\end{tcolorbox}
\section[\texorpdfstring{\sk{नीडम्} (nīḍam)}{नीडम् (nīḍam)}]{\sk{नीडम्} (nīḍam) — a nest}

\label{lesson:SA-C198-nidam}



\begin{quote}
Before the new one: say the Sanskrit for a leopard, then the Sanskrit for a rhinoceros.
\end{quote}

\subsection*{You'll want to know: \sk{नीडम्}}
\textbf{\sk{नीडम्}} — \emph{nīḍam} — \textquotedblleft{}a nest\textquotedblright{}.

\begin{grammarlens}[title={What you've built}]
One more word for this chapter.
\end{grammarlens}

\subsection*{Your turn}
\begin{itemize}
\raggedright
  \item \emph{Say it:} \emph{nīḍam}
  \item \emph{Say it:} \emph{nīḍam}, once more
  \item \emph{Say it:} say \emph{nīḍam}
  \item \emph{Recall:} say the Sanskrit for a leopard, then the Sanskrit for a rhinoceros, then say \emph{nīḍam} again
\end{itemize}

\begin{tcolorbox}[breakable,colback=teal!4,colframe=teal!35!black,title={Before you move on}]
What does \sk{नीडम्} mean? (\textquotedblleft{}A nest\textquotedblright{}.) Say it once more.
\end{tcolorbox}
\section[\texorpdfstring{\sk{अण्डम्} (aṇḍam)}{अण्डम् (aṇḍam)}]{\sk{अण्डम्} (aṇḍam) — an egg}

\label{lesson:SA-C198-andam}



\begin{quote}
Before the new one: say the Sanskrit for a rhinoceros, then the Sanskrit for a nest.
\end{quote}

\subsection*{You'll want to know: \sk{अण्डम्}}
\textbf{\sk{अण्डम्}} — \emph{aṇḍam} — \textquotedblleft{}an egg\textquotedblright{}.

\begin{grammarlens}[title={What you've built}]
One more word for this chapter.
\end{grammarlens}

\subsection*{Your turn}
\begin{itemize}
\raggedright
  \item \emph{Say it:} \emph{aṇḍam}
  \item \emph{Say it:} \emph{aṇḍam}, once more
  \item \emph{Say it:} say \emph{aṇḍam}
  \item \emph{Recall:} say the Sanskrit for a rhinoceros, then the Sanskrit for a nest, then say \emph{aṇḍam} again
\end{itemize}

\begin{tcolorbox}[breakable,colback=teal!4,colframe=teal!35!black,title={Before you move on}]
What does \sk{अण्डम्} mean? (\textquotedblleft{}An egg\textquotedblright{}.) Say it once more.
\end{tcolorbox}
